\documentclass[12pt,a4paper]{article}

% ── Packages ──────────────────────────────────────────────────────────────────
\usepackage{amsmath,amssymb}
\usepackage[margin=2.5cm]{geometry}
\usepackage{booktabs}
\usepackage{xcolor}
\usepackage{hyperref}
\usepackage{listings}
\usepackage{pgfplots}
\pgfplotsset{compat=1.18}
\usepackage[most,breakable]{tcolorbox}
\tcbuselibrary{skins,breakable}
\usepackage{enumitem}
\usepackage{fancyhdr}
\usepackage{tikz}
\usepackage{array}

% ── tcolorbox environments ─────────────────────────────────────────────────────
\newtcolorbox{curiositybox}[1]{colback=yellow!10,colframe=orange!80!black,fonttitle=\bfseries,title=#1,breakable}
\newtcolorbox{infobox}[1]{colback=blue!5,colframe=blue!60!black,fonttitle=\bfseries,title=#1,breakable}
\newtcolorbox{mistakebox}[1]{colback=red!5,colframe=red!60!black,fonttitle=\bfseries,title=#1,breakable}
\newtcolorbox{learnbox}[1]{colback=green!5,colframe=green!60!black,fonttitle=\bfseries,title=#1,breakable}

% ── lstlisting ─────────────────────────────────────────────────────────────────
\lstset{language=Python,basicstyle=\ttfamily\small,keywordstyle=\color{blue},
        commentstyle=\color{gray},stringstyle=\color{red},
        showstringspaces=false,breaklines=true,frame=single}

% ── fancyhdr ───────────────────────────────────────────────────────────────────
\pagestyle{fancy}
\fancyhf{}
\lhead{Topic 07: Harmonic Conjugate}
\rhead{Unit I --- Complex Variables}
\cfoot{\thepage}

% ── Title ──────────────────────────────────────────────────────────────────────
\title{\textbf{Topic 07: Finding Harmonic Conjugate}\\
       \large Unit I --- Complex Variable: Differentiation\\
       Complex Variables and Special Functions\\
       JNTUA College of Engineering (Autonomous)}
\author{B.Tech Engineering --- Lecture Notes}
\date{}

% ══════════════════════════════════════════════════════════════════════════════
\begin{document}
\maketitle
\tableofcontents
\newpage

% ══════════════════════════════════════════════════════════════════════════════
\section{Opening Hook}

\begin{curiositybox}{Hook}
Suppose someone gives you only the real part of an analytic function. Can you recover the
missing imaginary part? Surprisingly, yes --- if the given function is harmonic. The missing
partner is called the harmonic conjugate.
\end{curiositybox}

% ══════════════════════════════════════════════════════════════════════════════
\section{Why Harmonic Conjugates Matter}

In many engineering applications --- heat conduction, electrostatics, fluid flow --- the
governing equations reduce to Laplace's equation. When we model such a problem, we often
know one real-valued harmonic function (say the temperature distribution or the electric
potential) but need the companion function to form a complete analytic function. The harmonic
conjugate is precisely that companion.

\begin{center}
\begin{tabular}{@{}lll@{}}
\toprule
\textbf{Situation} & \textbf{Without conjugate} & \textbf{With conjugate} \\
\midrule
Heat flow & Only temperature $u$ known & Recover heat-flux potential $v$ \\
Electrostatics & Only potential $u$ known & Recover field lines via $v$ \\
Fluid mechanics & Only velocity potential $u$ & Recover stream function $v$ \\
Complex analysis & Isolated real part & Full analytic function $f(z)=u+iv$ \\
\bottomrule
\end{tabular}
\end{center}

\begin{learnbox}{Key Takeaway}
Whenever $u$ (or $v$) is harmonic, we can recover the partner function using the
Cauchy--Riemann (C-R) equations, giving us a full analytic function $f(z)=u+iv$.
\end{learnbox}

% ══════════════════════════════════════════════════════════════════════════════
\section{Definition}

\begin{infobox}{Harmonic Conjugate --- Formal Definition}
Let $u(x,y)$ be harmonic in a simply connected domain $D$ (i.e.\
$u_{xx}+u_{yy}=0$ throughout $D$). A function $v(x,y)$ is called the
\textbf{harmonic conjugate} of $u$ if
\[
  \frac{\partial u}{\partial x}=\frac{\partial v}{\partial y}
  \qquad\text{and}\qquad
  \frac{\partial u}{\partial y}=-\frac{\partial v}{\partial x}
\]
(the Cauchy--Riemann equations), so that $f(z)=u(x,y)+iv(x,y)$ is analytic in $D$.

\medskip
\textbf{Notes:}
\begin{itemize}[leftmargin=1.5em]
  \item $v$ is itself harmonic once $u$ is harmonic and the C-R equations hold.
  \item The harmonic conjugate is unique \emph{up to an additive real constant}.
  \item $v$ is \emph{not} the harmonic conjugate of $u$ --- the relationship is
        ordered: $v$ is the conjugate of $u$, not vice versa (though $-u$ is the
        harmonic conjugate of $v$).
\end{itemize}
\end{infobox}

% ══════════════════════════════════════════════════════════════════════════════
\section{Method Using C-R Equations}

\begin{infobox}{Step-by-Step Procedure (given $u$, find $v$)}
\textbf{Step 1.} Verify (or assume) $u$ satisfies $u_{xx}+u_{yy}=0$.

\textbf{Step 2.} Compute $u_x = \dfrac{\partial u}{\partial x}$ and
                               $u_y = \dfrac{\partial u}{\partial y}$.

\textbf{Step 3.} Use $v_y = u_x$ to write $v = \displaystyle\int u_x\,dy + \phi(x)$,
where $\phi(x)$ is an \emph{unknown function of $x$ alone} (the ``constant'' of integration
in $y$).

\textbf{Step 4.} Differentiate the result with respect to $x$:
$v_x = \dfrac{\partial}{\partial x}\!\left[\int u_x\,dy\right] + \phi'(x)$.

\textbf{Step 5.} Apply the second C-R equation: $v_x = -u_y$.
Set the expression from Step 4 equal to $-u_y$ and solve for $\phi'(x)$.

\textbf{Step 6.} Integrate to find $\phi(x)$ (a constant $c$ appears here).

\textbf{Step 7.} Write the complete $v(x,y)$ and, if required, $f(z)=u+iv$.
\end{infobox}

% ══════════════════════════════════════════════════════════════════════════════
\section{Alternative Starting from $v$}

\begin{infobox}{Step-by-Step Procedure (given $v$, find $u$)}
\textbf{Step 1.} Verify $v_{xx}+v_{yy}=0$.

\textbf{Step 2.} Compute $v_x$ and $v_y$.

\textbf{Step 3.} Use $u_x = v_y$, so $u = \displaystyle\int v_y\,dx + \psi(y)$,
where $\psi(y)$ is an unknown function of $y$ alone.

\textbf{Step 4.} Differentiate: $u_y = \dfrac{\partial}{\partial y}\!\left[\int v_y\,dx\right] + \psi'(y)$.

\textbf{Step 5.} Apply $u_y = -v_x$. Solve for $\psi'(y)$ and integrate.

\textbf{Step 6.} Write $u(x,y)$ and $f(z)=u+iv$.
\end{infobox}

\begin{learnbox}{Remember}
Whether you start from $u$ or $v$, the key is: one C-R equation gives the integration
direction; the other determines the unknown function of the remaining variable.
\end{learnbox}

% ══════════════════════════════════════════════════════════════════════════════
\section{Uniqueness Up to an Additive Constant}

\begin{infobox}{Uniqueness Theorem}
If $v_1$ and $v_2$ are both harmonic conjugates of the same $u$ in a connected domain $D$,
then $v_1 - v_2 = c$ for some real constant $c$.
\end{infobox}

\textbf{Proof sketch.} Since both $v_1$ and $v_2$ satisfy the C-R equations with the same
$u$, the function $w = v_1 - v_2$ satisfies $(w)_x = 0$ and $(w)_y = 0$ throughout $D$.
In a connected open set, a function with zero gradient is constant. Hence $v_1-v_2=c$.

\medskip
\textbf{Practical consequence.} When we integrate to find $v$, the arbitrary real constant
$c$ is unavoidable and acceptable --- it represents the freedom to shift the level curves
of $v$ without changing the analytic function's analyticity.

\begin{learnbox}{Uniqueness Takeaway}
Every harmonic conjugate of $u$ in a connected domain $D$ differs from every other harmonic
conjugate of $u$ by at most a real constant. We write $v+c$, $c\in\mathbb{R}$.
\end{learnbox}

% ══════════════════════════════════════════════════════════════════════════════
\section{Worked Examples}

% ── Example 1 ─────────────────────────────────────────────────────────────────
\subsection*{Example 1: $u = x^2 - y^2$}

\textbf{Given:} $u(x,y)=x^2-y^2$. Find the harmonic conjugate $v$.

\medskip
\textbf{Step 1 --- Check harmonicity:}
\[
  u_{xx} = 2,\quad u_{yy}=-2,\quad u_{xx}+u_{yy}=0.\quad\checkmark
\]

\textbf{Step 2 --- Compute partial derivatives:}
\[
  u_x = 2x,\qquad u_y = -2y.
\]

\textbf{Step 3 --- Integrate $u_x$ w.r.t.\ $y$:}
\[
  v = \int 2x\,dy + \phi(x) = 2xy + \phi(x).
\]

\textbf{Step 4 --- Differentiate w.r.t.\ $x$:}
\[
  v_x = 2y + \phi'(x).
\]

\textbf{Step 5 --- Apply $v_x = -u_y$:}
\[
  2y+\phi'(x)=-(-2y)=2y \implies \phi'(x)=0 \implies \phi(x)=c.
\]

\textbf{Result:}
\[
  \boxed{v(x,y)=2xy+c,\qquad f(z)=(x^2-y^2)+i(2xy)+ic=z^2+ic.}
\]

\begin{learnbox}{Example 1 Key Result}
$u=x^2-y^2$ leads to $v=2xy+c$. The analytic function is $f(z)=z^2+ic$. This is the
classic example showing the real and imaginary parts of $z^2$.
\end{learnbox}

% ── Example 2 ─────────────────────────────────────────────────────────────────
\subsection*{Example 2: $u = e^x\cos y$}

\textbf{Given:} $u(x,y)=e^x\cos y$. Find $v$ and $f(z)$.

\medskip
\textbf{Step 1:}
\[
  u_{xx}=e^x\cos y,\quad u_{yy}=-e^x\cos y,\quad u_{xx}+u_{yy}=0.\quad\checkmark
\]

\textbf{Step 2:}
\[
  u_x = e^x\cos y,\qquad u_y = -e^x\sin y.
\]

\textbf{Step 3:}
\[
  v = \int e^x\cos y\,dy + \phi(x) = e^x\sin y + \phi(x).
\]

\textbf{Step 4:}
\[
  v_x = e^x\sin y + \phi'(x).
\]

\textbf{Step 5:}
\[
  e^x\sin y+\phi'(x)= -(-e^x\sin y)=e^x\sin y \implies \phi'(x)=0 \implies \phi(x)=c.
\]

\textbf{Result:}
\[
  \boxed{v=e^x\sin y+c,\qquad f(z)=e^x(\cos y+i\sin y)+ic=e^z+ic.}
\]

\begin{learnbox}{Example 2 Key Result}
$u=e^x\cos y$ gives $v=e^x\sin y+c$ and $f(z)=e^z+ic$. Real and imaginary parts of
$e^z$ are $e^x\cos y$ and $e^x\sin y$.
\end{learnbox}

% ── Example 3 ─────────────────────────────────────────────────────────────────
\subsection*{Example 3: $u = \dfrac{x}{x^2+y^2}$}

\textbf{Given:} $u=\dfrac{x}{x^2+y^2}$ (defined for $z\neq0$). Find $v$.

\medskip
\textbf{Step 1 --- Check harmonicity:} Let $r^2=x^2+y^2$.
\[
  u_x = \frac{(x^2+y^2)-x\cdot 2x}{(x^2+y^2)^2}=\frac{y^2-x^2}{r^4},\qquad
  u_y = \frac{-2xy}{r^4}.
\]
\[
  u_{xx}=\frac{-2x(y^2-x^2)\cdot2/r^4 - (y^2-x^2)\cdot 4x/r^4}{1}
  \quad\text{(detailed computation yields }u_{xx}+u_{yy}=0\text{).}\quad\checkmark
\]

\textbf{Step 3:}
\[
  v = \int \frac{y^2-x^2}{(x^2+y^2)^2}\,dy + \phi(x).
\]
Use the standard integral result for this form:
\[
  \int \frac{y^2-x^2}{(x^2+y^2)^2}\,dy = -\frac{y}{x^2+y^2}+\phi(x).
\]

\textbf{Step 4--5:}
\[
  v_x = \frac{2xy}{(x^2+y^2)^2}+\phi'(x).
\]
\[
  -u_y = \frac{2xy}{(x^2+y^2)^2} \implies \phi'(x)=0 \implies \phi(x)=c.
\]

\textbf{Result:}
\[
  \boxed{v=-\frac{y}{x^2+y^2}+c.}
\]
Note: $u+iv = \dfrac{x-iy}{x^2+y^2}=\dfrac{\overline{z}}{|z|^2}=\dfrac{1}{z}+ic$,
confirming $f(z)=1/z+ic$.

\begin{learnbox}{Example 3 Key Result}
$u=x/(x^2+y^2)$ gives $v=-y/(x^2+y^2)+c$ and $f(z)=1/z+ic$. This is a rational
function with a singularity at the origin.
\end{learnbox}

% ── Example 4 ─────────────────────────────────────────────────────────────────
\subsection*{Example 4: $u = x^3 - 3xy^2$}

\textbf{Given:} $u=x^3-3xy^2$. Find $v$ and $f(z)$.

\medskip
\textbf{Step 1:}
\[
  u_{xx}=6x,\quad u_{yy}=-6x,\quad u_{xx}+u_{yy}=0.\quad\checkmark
\]

\textbf{Step 2:}
\[
  u_x=3x^2-3y^2,\qquad u_y=-6xy.
\]

\textbf{Step 3:}
\[
  v=\int(3x^2-3y^2)\,dy+\phi(x)=3x^2y-y^3+\phi(x).
\]

\textbf{Step 4--5:}
\[
  v_x=6xy+\phi'(x)=-u_y=6xy \implies \phi'(x)=0 \implies \phi(x)=c.
\]

\textbf{Result:}
\[
  \boxed{v=3x^2y-y^3+c,\qquad f(z)=(x^3-3xy^2)+i(3x^2y-y^3)+ic=z^3+ic.}
\]

\begin{learnbox}{Example 4 Key Result}
$u=x^3-3xy^2$ (real part of $z^3$) gives $v=3x^2y-y^3+c$.
\end{learnbox}

% ── Example 5 ─────────────────────────────────────────────────────────────────
\subsection*{Example 5: Given $v = e^x\sin y$, find $u$}

\textbf{Given:} $v(x,y)=e^x\sin y$. Find $u$.

\medskip
\textbf{Step 1:}
\[
  v_{xx}=e^x\sin y,\quad v_{yy}=-e^x\sin y,\quad v_{xx}+v_{yy}=0.\quad\checkmark
\]

\textbf{Step 2:}
\[
  v_x=e^x\sin y,\qquad v_y=e^x\cos y.
\]

\textbf{Step 3 --- Use $u_x=v_y$:}
\[
  u=\int e^x\cos y\,dx+\psi(y)=e^x\cos y+\psi(y).
\]

\textbf{Step 4--5 --- Use $u_y=-v_x$:}
\[
  u_y=-e^x\sin y+\psi'(y)=-v_x=-e^x\sin y \implies \psi'(y)=0 \implies \psi(y)=c.
\]

\textbf{Result:}
\[
  \boxed{u=e^x\cos y+c,\qquad f(z)=e^z+c.}
\]

\begin{learnbox}{Example 5 Key Result}
Starting from $v=e^x\sin y$ recovers $u=e^x\cos y+c$, confirming $f(z)=e^z+c$.
\end{learnbox}

% ── Example 6 ─────────────────────────────────────────────────────────────────
\subsection*{Example 6: Checking Harmonicity First --- $u = x^2 + y^2$}

\textbf{Given:} $u=x^2+y^2$. Can a harmonic conjugate exist?

\medskip
\textbf{Step 1:}
\[
  u_{xx}=2,\quad u_{yy}=2,\quad u_{xx}+u_{yy}=4\neq0.
\]

Since $u$ is \textbf{not harmonic}, no harmonic conjugate exists.

\begin{learnbox}{Example 6 Key Result}
Always check $\nabla^2 u=0$ first. If the check fails, stop immediately --- there is no
analytic $f(z)=u+iv$ for this $u$.
\end{learnbox}

% ── Example 7 ─────────────────────────────────────────────────────────────────
\subsection*{Example 7: $u = \sin x\cosh y$}

\textbf{Given:} $u=\sin x\cosh y$. Find $v$.

\medskip
\textbf{Step 1:}
\[
  u_{xx}=-\sin x\cosh y,\quad u_{yy}=\sin x\cosh y,\quad u_{xx}+u_{yy}=0.\quad\checkmark
\]

\textbf{Step 2:}
\[
  u_x=\cos x\cosh y,\qquad u_y=\sin x\sinh y.
\]

\textbf{Step 3:}
\[
  v=\int\cos x\cosh y\,dy+\phi(x)=\cos x\sinh y+\phi(x).
\]

\textbf{Step 4--5:}
\[
  v_x=-\sin x\sinh y+\phi'(x)=-u_y=-\sin x\sinh y \implies \phi'(x)=0 \implies \phi(x)=c.
\]

\textbf{Result:}
\[
  \boxed{v=\cos x\sinh y+c,\qquad f(z)=\sin z+c.}
\]

\begin{learnbox}{Example 7 Key Result}
$u=\sin x\cosh y$, $v=\cos x\sinh y+c$, $f(z)=\sin z+c$. The hyperbolic functions arise
naturally from the trigonometric functions of a complex argument.
\end{learnbox}

% ── Example 8 ─────────────────────────────────────────────────────────────────
\subsection*{Example 8: $u = \ln\!\sqrt{x^2+y^2}$}

\textbf{Given:} $u=\ln\sqrt{x^2+y^2}=\tfrac{1}{2}\ln(x^2+y^2)$ for $z\neq0$. Find $v$.

\medskip
\textbf{Step 1:}
\[
  u_x=\frac{x}{x^2+y^2},\quad u_y=\frac{y}{x^2+y^2}.
\]
\[
  u_{xx}=\frac{y^2-x^2}{(x^2+y^2)^2},\quad u_{yy}=\frac{x^2-y^2}{(x^2+y^2)^2},
  \quad u_{xx}+u_{yy}=0.\quad\checkmark
\]

\textbf{Step 3:}
\[
  v=\int\frac{x}{x^2+y^2}\,dy+\phi(x)=\arctan\!\left(\frac{y}{x}\right)+\phi(x).
\]

\textbf{Step 4--5:}
\[
  v_x=\frac{-y/x^2}{1+(y/x)^2}+\phi'(x)=\frac{-y}{x^2+y^2}+\phi'(x).
\]
\[
  -u_y=-\frac{y}{x^2+y^2} \implies \phi'(x)=0 \implies \phi(x)=c.
\]

\textbf{Result:}
\[
  \boxed{v=\arctan\!\left(\frac{y}{x}\right)+c=\arg(z)+c,
  \qquad f(z)=\ln|z|+i\arg(z)+ic=\operatorname{Log}z+ic.}
\]

\begin{learnbox}{Example 8 Key Result}
$u=\ln|z|$ gives $v=\arg(z)+c$. The analytic function is the principal logarithm
$\operatorname{Log}z$. This connects the geometric notion of argument to the harmonic
conjugate.
\end{learnbox}

% ══════════════════════════════════════════════════════════════════════════════
\section{Procedure Summary Table (MANDATORY UPGRADE)}

\begin{center}
\begin{tabular}{@{}>{\bfseries}c p{5cm} p{6cm}@{}}
\toprule
Step & What to Do & Why It Is Needed \\
\midrule
1 & Verify $u_{xx}+u_{yy}=0$ &
    Ensures a harmonic conjugate actually exists; skip if given, verify if not stated \\
\midrule
2 & Compute $u_x$ and $u_y$ &
    These feed directly into the two C-R equations \\
\midrule
3 & Set $v_y=u_x$; integrate w.r.t.\ $y$ &
    One C-R equation gives $v$ as an integral in $y$; the ``constant'' is $\phi(x)$ \\
\midrule
4 & Differentiate result w.r.t.\ $x$ &
    Prepares $v_x$ for comparison with the second C-R equation \\
\midrule
5 & Apply $v_x=-u_y$; solve for $\phi'(x)$ &
    The second C-R equation determines the unknown function $\phi(x)$ \\
\midrule
6 & Integrate $\phi'(x)$ to get $\phi(x)=c$ (or function) &
    Recovers $\phi(x)$ fully; $c$ is the additive freedom \\
\midrule
7 & Write $v=(\text{result from Step 3})+\phi(x)$ &
    Assembles the final harmonic conjugate \\
\midrule
8 & (Optional) Form $f(z)=u+iv$ and express in terms of $z$ &
    Confirms correctness and gives the analytic function in closed form \\
\bottomrule
\end{tabular}
\end{center}

% ══════════════════════════════════════════════════════════════════════════════
\section{Excel Example (MANDATORY)}

\textbf{Scenario:} Numerically verify that $u=x^2-y^2$ and $v=2xy$ satisfy the C-R
equations by computing finite-difference approximations of the partial derivatives.

\medskip
\textbf{Spreadsheet layout} (start in cell A1):

\begin{center}
\begin{tabular}{@{}ccccccccc@{}}
\toprule
Col A & Col B & Col C & Col D & Col E & Col F & Col G & Col H \\
$x$ & $y$ & $u=A^2-B^2$ & $v=2AB$ & $u_x\approx$ & $u_y\approx$ & $v_x\approx$ & $v_y\approx$ \\
\midrule
1.0 & 1.0 & \texttt{=A2\^{}2-B2\^{}2} & \texttt{=2*A2*B2} & \texttt{=2*A2} & \texttt{=-2*B2} & \texttt{=2*B2} & \texttt{=2*A2} \\
2.0 & 0.5 & \texttt{=A3\^{}2-B3\^{}2} & \texttt{=2*A3*B3} & \texttt{=2*A3} & \texttt{=-2*B3} & \texttt{=2*B3} & \texttt{=2*A3} \\
3.0 & 2.0 & \texttt{=A4\^{}2-B4\^{}2} & \texttt{=2*A4*B4} & \texttt{=2*A4} & \texttt{=-2*B4} & \texttt{=2*B4} & \texttt{=2*A4} \\
\bottomrule
\end{tabular}
\end{center}

\textbf{Verification columns to add:}

\begin{center}
\begin{tabular}{@{}cll@{}}
\toprule
Column & Header & Formula \\
\midrule
Col I & $u_x - v_y$ (should be 0) & \texttt{=E2-H2} \\
Col J & $u_y + v_x$ (should be 0) & \texttt{=F2+G2} \\
\bottomrule
\end{tabular}
\end{center}

\textbf{Sample computed values:}

\begin{center}
\begin{tabular}{@{}cccccccc@{}}
\toprule
$x$ & $y$ & $u$ & $v$ & $u_x$ & $u_y$ & $v_x$ & $v_y$ \\
\midrule
1.0 & 1.0 & 0.00 & 2.00 & 2.00 & $-2.00$ & 2.00 & 2.00 \\
2.0 & 0.5 & 3.75 & 2.00 & 4.00 & $-1.00$ & 1.00 & 4.00 \\
3.0 & 2.0 & 5.00 & 12.00 & 6.00 & $-4.00$ & 4.00 & 6.00 \\
\bottomrule
\end{tabular}
\end{center}

\textbf{Observation:} Column I ($u_x-v_y$) and Column J ($u_y+v_x$) are identically
zero for all rows, confirming the C-R equations hold numerically.

% ══════════════════════════════════════════════════════════════════════════════
\section{Python Example (MANDATORY)}

\textbf{Task:} Use SymPy to (1) check harmonicity of a given $u$, (2) derive the harmonic
conjugate $v$ symbolically, and (3) form $f(z)=u+iv$.

\begin{lstlisting}[language=Python,caption={Harmonic Conjugate via SymPy}]
import sympy as sp

x, y, c = sp.symbols('x y c', real=True)

# ── Example: u = x^2 - y^2 ──────────────────────────────────────────────────
u = x**2 - y**2

# Step 1: Check harmonicity
laplacian_u = sp.diff(u, x, 2) + sp.diff(u, y, 2)
print("Laplacian of u:", laplacian_u)  # Expected: 0

# Step 2: Partial derivatives
u_x = sp.diff(u, x)
u_y = sp.diff(u, y)
print("u_x =", u_x)  # 2*x
print("u_y =", u_y)  # -2*y

# Step 3: Integrate u_x w.r.t. y  =>  v = integral(u_x, y) + phi(x)
phi = sp.Function('phi')
v_partial = sp.integrate(u_x, y)
print("v (partial) =", v_partial)  # 2*x*y

# Step 4: Differentiate w.r.t. x and apply v_x = -u_y
# v_x = diff(v_partial, x) + phi'(x) = -u_y
phi_x = sp.Symbol('phi_x')  # represents phi'(x)
eq = sp.Eq(sp.diff(v_partial, x) + phi_x, -u_y)
print("Equation for phi'(x):", eq)   # phi_x = 0

phi_prime = sp.solve(eq, phi_x)[0]
print("phi'(x) =", phi_prime)  # 0

# Step 5: Integrate phi'(x)
phi_val = sp.integrate(phi_prime, x)
print("phi(x) =", phi_val, "+ c")  # 0 + c

# Step 6: Complete harmonic conjugate
v = v_partial + phi_val + c
print("v(x,y) =", v)  # 2*x*y + c

# Step 7: Analytic function
z = x + sp.I * y
f_z = u + sp.I * v
f_simplified = sp.expand(f_z)
print("f(z) = u + iv =", f_simplified)  # x**2 - y**2 + 2*I*x*y + I*c

# ── Bonus: verify with a second example u = e^x * cos(y) ───────────────────
u2 = sp.exp(x) * sp.cos(y)
lap2 = sp.diff(u2, x, 2) + sp.diff(u2, y, 2)
print("\nLaplacian of e^x*cos(y):", sp.simplify(lap2))  # 0

u2_x = sp.diff(u2, x)
v2_partial = sp.integrate(u2_x, y)
print("v2 (partial) =", v2_partial)   # exp(x)*sin(y)

phi2_prime = sp.solve(sp.Eq(sp.diff(v2_partial, x) + phi_x, -sp.diff(u2, y)), phi_x)[0]
print("phi2'(x) =", phi2_prime)  # 0

v2 = v2_partial + phi2_prime + c
print("v2 =", v2)   # exp(x)*sin(y) + c
print("f2(z) = e^z + c (analytic function)")
\end{lstlisting}

\textbf{Expected output (summary):}
\begin{verbatim}
Laplacian of u: 0
u_x = 2*x
u_y = -2*y
v (partial) = 2*x*y
phi'(x) = 0
phi(x) = 0 + c
v(x,y) = 2*x*y + c
f(z) = u + iv = x**2 - y**2 + 2*I*x*y + I*c

Laplacian of e^x*cos(y): 0
v2 (partial) = exp(x)*sin(y)
phi2'(x) = 0
v2 = exp(x)*sin(y) + c
f2(z) = e^z + c (analytic function)
\end{verbatim}

% ══════════════════════════════════════════════════════════════════════════════
\section{Viva Questions}

\begin{enumerate}[label=\textbf{Q\arabic*.}]

\item \textbf{What is the harmonic conjugate of a function $u$?}\\
      The harmonic conjugate of $u$ is a function $v$ such that $v_y=u_x$ and $v_x=-u_y$
      (the Cauchy--Riemann equations), making $f(z)=u+iv$ analytic.

\item \textbf{Is a harmonic conjugate always unique?}\\
      No. The harmonic conjugate is unique \emph{only up to an additive real constant}.
      If $v$ is a harmonic conjugate of $u$, so is $v+c$ for any $c\in\mathbb{R}$.

\item \textbf{What must you check before finding a harmonic conjugate?}\\
      You must verify that $u$ satisfies Laplace's equation: $u_{xx}+u_{yy}=0$ in the
      domain. If $u$ is not harmonic, no harmonic conjugate exists.

\item \textbf{In the method using C-R equations, why does an unknown function $\phi(x)$
      appear rather than a plain constant?}\\
      Because we integrate $v_y=u_x$ with respect to $y$, the ``constant'' of integration
      can depend on the remaining variable $x$. It is determined by applying the second
      C-R equation $v_x=-u_y$.

\item \textbf{If $v$ is the harmonic conjugate of $u$, is $u$ the harmonic conjugate of
      $v$?}\\
      No. If $f=u+iv$ is analytic, then $-u$ is the harmonic conjugate of $v$. The roles
      of $u$ and $v$ are not symmetric.

\item \textbf{What is the harmonic conjugate of $u=\ln|z|$?}\\
      The harmonic conjugate is $v=\arg(z)+c$, so $f(z)=\operatorname{Log}z+ic$.

\item \textbf{Does a harmonic function always have a harmonic conjugate?}\\
      In a simply connected domain, yes. In a multiply connected domain (e.g., a
      punctured disc), a harmonic conjugate may not exist globally (e.g., $u=\ln|z|$
      has no single-valued harmonic conjugate around a loop encircling the origin).

\item \textbf{What is the geometric relationship between the level curves of $u$ and $v$?}\\
      The level curves of $u$ (equipotentials) and the level curves of $v$ (stream lines)
      are orthogonal families of curves wherever $\nabla u\neq\mathbf{0}$.

\end{enumerate}

% ══════════════════════════════════════════════════════════════════════════════
\section{Descriptive Questions}

\begin{enumerate}[label=\textbf{D\arabic*.}]

\item Define harmonic conjugate. State and prove the uniqueness theorem for harmonic
      conjugates in a connected domain. Illustrate with the pair $u=x^2-y^2$,
      $v=2xy+c$.

\item Derive the step-by-step procedure for finding the harmonic conjugate of a given
      harmonic function $u(x,y)$ using the Cauchy--Riemann equations. Apply the method
      to find the harmonic conjugate of $u=e^x\cos y$ and write down the corresponding
      analytic function $f(z)$.

\item Explain how to find $u$ when $v$ is given. Use $v=2xy+3x$ as an example: verify
      harmonicity, find $u$, and write $f(z)$.

\item Find the harmonic conjugate of $u=\ln\sqrt{x^2+y^2}$ for $z\neq0$. Identify the
      resulting analytic function and discuss the domain of validity.

\item A student claims that $u=x^2+y^2$ has a harmonic conjugate. Identify the error
      in this claim, and explain what additional conditions on $u$ are necessary for a
      harmonic conjugate to exist. Give two examples of functions that do satisfy these
      conditions.

\end{enumerate}

% ══════════════════════════════════════════════════════════════════════════════
\section{Practice Problems}

\begin{enumerate}[label=\textbf{P\arabic*.}]

\item Find the harmonic conjugate of $u=2xy+x$.\\
      \textit{Hint:} $u_x=2y+1$; integrate w.r.t.\ $y$; expected $v=y^2-x^2+y+c$.

\item Given $v=x^3y-xy^3+y$, verify $v$ is harmonic and find $u$.\\
      \textit{Hint:} Use $u_x=v_y$; expected $u=x^2y^2-\tfrac{1}{2}(x^4+y^4)/2+\cdots$;
      work step by step.

\item Find the harmonic conjugate of $u=\cos x\cosh y$.\\
      \textit{Hint:} $u_x=-\sin x\cosh y$; integrate w.r.t.\ $y$; expected
      $v=-\sin x\sinh y+c$; $f(z)=\cos z+c$.

\item Show that $u=x/(x^2+y^2)$ is harmonic and find $v$.\\
      \textit{Hint:} $v=-y/(x^2+y^2)+c$; $f(z)=1/z+ic$.

\item Given $u=e^{2x}\cos 2y$, find the harmonic conjugate and the analytic function.\\
      \textit{Hint:} $v=e^{2x}\sin 2y+c$; $f(z)=e^{2z}+c$.

\item Verify that $u=x^4-6x^2y^2+y^4$ is harmonic and find its harmonic conjugate.\\
      \textit{Hint:} $u=\operatorname{Re}(z^4)$; expected $v=4x^3y-4xy^3+c=\operatorname{Im}(z^4)+c$.

\end{enumerate}

% ══════════════════════════════════════════════════════════════════════════════
\section{MCQs}

\begin{enumerate}[label=\textbf{MCQ\arabic*.}]

\item If $u=x^2-y^2$, the harmonic conjugate $v$ is:
      \begin{enumerate}[label=(\alph*)]
        \item $x^2+y^2+c$
        \item $-2xy+c$
        \item \textbf{$2xy+c$}
        \item $x^2-y^2+c$
      \end{enumerate}
      \textit{Explanation:} $v_y=u_x=2x \Rightarrow v=2xy+\phi(x)$; $\phi'(x)=0$ gives
      $v=2xy+c$.

\item The harmonic conjugate is unique:
      \begin{enumerate}[label=(\alph*)]
        \item Always, with no freedom
        \item Up to a multiplicative constant
        \item \textbf{Up to an additive real constant}
        \item Only if the domain is multiply connected
      \end{enumerate}
      \textit{Explanation:} Any two harmonic conjugates of $u$ differ by a real constant,
      as shown by the uniqueness theorem.

\item If $u$ is harmonic in a simply connected domain, which of the following must hold
      for $v$ to be its harmonic conjugate?
      \begin{enumerate}[label=(\alph*)]
        \item $v_x=u_x$ and $v_y=u_y$
        \item $v_x=u_y$ and $v_y=u_x$
        \item \textbf{$v_x=-u_y$ and $v_y=u_x$}
        \item $v_x=u_y$ and $v_y=-u_x$
      \end{enumerate}
      \textit{Explanation:} These are precisely the Cauchy--Riemann equations.

\item The harmonic conjugate of $u=e^x\cos y$ is:
      \begin{enumerate}[label=(\alph*)]
        \item $e^x\cos y+c$
        \item $-e^x\sin y+c$
        \item $e^y\sin x+c$
        \item \textbf{$e^x\sin y+c$}
      \end{enumerate}
      \textit{Explanation:} Integrating $v_y=u_x=e^x\cos y$ with respect to $y$ gives
      $v=e^x\sin y+\phi(x)$; the second C-R equation forces $\phi'(x)=0$.

\item Which of the following functions does \textbf{not} have a harmonic conjugate?
      \begin{enumerate}[label=(\alph*)]
        \item $u=\ln\sqrt{x^2+y^2}$
        \item $u=x^2-y^2$
        \item \textbf{$u=x^2+y^2$}
        \item $u=e^x\cos y$
      \end{enumerate}
      \textit{Explanation:} $\nabla^2(x^2+y^2)=4\neq0$, so it is not harmonic and has no
      harmonic conjugate.

\end{enumerate}

% ══════════════════════════════════════════════════════════════════════════════
\section{Common Mistakes}

\begin{mistakebox}{Frequent Errors When Finding Harmonic Conjugates}
\begin{tabular}{@{}>{\bfseries}p{4cm} p{5cm} p{5cm}@{}}
\toprule
Mistake & Example of the Error & Correct Approach \\
\midrule
Forgetting the arbitrary constant &
  Writing $v=2xy$ for $u=x^2-y^2$ &
  Always append $+c$: $v=2xy+c$ \\
\midrule
Using a plain constant instead of $\phi(x)$ &
  Writing $v=2xy+A$ without deriving $A$ &
  Use $\phi(x)$ (function of $x$); determine by the second C-R equation \\
\midrule
Sign error in C-R equation &
  Writing $v_x=u_y$ instead of $v_x=-u_y$ &
  The second C-R equation is $v_x=-u_y$; the minus sign is essential \\
\midrule
Not checking harmonicity first &
  Attempting to find conjugate of $u=x^2+y^2$ &
  Always verify $u_{xx}+u_{yy}=0$ before starting \\
\midrule
Integrating w.r.t.\ wrong variable &
  Integrating $u_x$ w.r.t.\ $x$ instead of $y$ &
  $v_y=u_x$ means integrate $u_x$ w.r.t.\ $y$ \\
\midrule
Neglecting that $\phi$ may be non-constant &
  Assuming $\phi'(x)=0$ before computing it &
  Compute $v_x$, set equal to $-u_y$, and \emph{then} solve for $\phi'(x)$ \\
\bottomrule
\end{tabular}
\end{mistakebox}

% ══════════════════════════════════════════════════════════════════════════════
\section{Quick Recap}

\begin{learnbox}{Quick Recap --- Topic 07: Harmonic Conjugate}
\begin{itemize}[leftmargin=1.5em]
  \item $v$ is the harmonic conjugate of $u$ if and only if $v_y=u_x$ and $v_x=-u_y$
        (C-R equations hold).
  \item Prerequisite: $u_{xx}+u_{yy}=0$; if $\nabla^2u\neq0$, no conjugate exists.
  \item Finding $v$ from $u$: integrate $v_y=u_x$ w.r.t.\ $y$, introduce $\phi(x)$,
        apply $v_x=-u_y$ to find $\phi'(x)$, integrate once more.
  \item Finding $u$ from $v$: integrate $u_x=v_y$ w.r.t.\ $x$, introduce $\psi(y)$,
        apply $u_y=-v_x$ to find $\psi'(y)$.
  \item Uniqueness: the harmonic conjugate is unique up to an additive real constant $c$.
  \item Key pairs: $z^n$, $e^z$, $\sin z$, $\cos z$, $\operatorname{Log}z$, $1/z$ each
        give textbook examples of the $u$--$v$ relationship.
  \item Level curves of $u$ and $v$ are orthogonal wherever $\nabla u\neq\mathbf{0}$.
  \item Every worked solution must include the arbitrary constant --- omitting it is a
        standard examination error.
\end{itemize}
\end{learnbox}

\end{document}
